\newpage
\section{Related Work}

APEIRON-Core occupies a unique intersection of hardware abstraction, high-performance computing, and autonomous systems. We survey relevant prior work along three axes: deep learning frameworks, FPGA toolchains, and quantum computing platforms.

\subsection{Deep Learning Frameworks}
PyTorch~\cite{paszke2019pytorch} and TensorFlow~\cite{abadi2016tensorflow} provide robust GPU acceleration via CUDA and cuDNN, with recent extensions supporting Apple Silicon and AMD GPUs. However, their abstraction is tightly coupled to tensor operations and neural network layers. They do not natively support the continuous field semantics required by the APEIRON substrate (Layers~1–4), nor do they offer built-in fallback to FPGA or quantum backends. ONNX Runtime~\cite{onnxruntime2024} provides a more general execution provider model, but it is optimized for inference graphs rather than the open-ended field evolution central to APEIRON.

\subsection{FPGA Development}
Xilinx's PYNQ~\cite{pynq2024} framework enables Python-based overlays on Zynq devices, significantly lowering the barrier to FPGA acceleration. FINN~\cite{blott2018finn} and hls4ml~\cite{duarte2018hls4ml} generate dataflow architectures for neural networks. These tools, however, are specialized for deep learning inference. APEIRON-Core abstracts the FPGA as a general computational fabric for continuous field operations, with a custom bitstream that implements parallel field updates and interference detection. The fixed-point conversion routines address the impedance mismatch between Python's floating-point and AXI register interfaces.

\subsection{Quantum Computing Platforms}
Qiskit~\cite{qiskit2023} and Cirq~\cite{cirq2024} provide comprehensive Python APIs for quantum circuit construction and simulation. They support multiple backends, from local simulators to IBM Quantum hardware. The quantum backend in APEIRON-Core leverages Qiskit's AerSimulator and StatevectorSimulator but adds two critical enhancements: (1) a corrected SWAP-test implementation that correctly computes state overlap by examining the ancilla qubit, and (2) a partial trace implementation via NumPy reshaping that avoids the overhead of multiple circuit evaluations. These improvements are essential for the efficient, repeated coherence measurements required by APEIRON's Layer~4 dynamics.

\subsection{Hardware Abstraction in Autonomous Systems}
The Robot Operating System (ROS)~\cite{ros2024} provides a hardware abstraction layer for sensors and actuators but does not address computational heterogeneity. Similarly, the OpenCL~\cite{opencl2024} and SYCL~\cite{sycl2024} standards offer portable parallelism across CPUs, GPUs, and FPGAs, yet they lack the higher-level semantics of continuous fields and the integrated error-handling and cost-accounting features of APEIRON-Core. To our knowledge, no existing framework unifies CPU, GPU, FPGA, and quantum backends under a single API with the specific goal of supporting non-anthropocentric knowledge genesis.